% !TeX root = ../概率论与数理统计笔记.tex
% 由主文档通过 \input 载入；请编译主文档。
\setcounter{chapter}{4} % 保留原书章号。
\chapter{数学期望与数字特征}
\label{chap:ps-moments}

本章研究怎样用少量数值描述分布，以及这些数值遗漏了什么信息。
先修内容为单变量与联合分布；学习目标是正确计算期望、方差和协方差，并在计算前检查存在条件。

\section{期望的定义与存在性}
\label{sec:ps-expectation-definition}
定义\emph{数学期望}（mathematical expectation），由离散加权和与密度积分建立统一认识。
区分非负随机变量的扩展期望、一般随机变量的期望存在与有限可积；有限期望采用 $\mathbb E[|X|]<\infty$ 作为条件。
利用重尾例子说明对称性不能替代可积性，不能用发散正负项的抵消定义期望。

\section{期望的计算与基本性质}
\label{sec:ps-expectation-properties}
建立随机变量函数的期望公式、线性性、单调性以及指示变量方法。
说明有限和的期望线性性不要求独立，而乘积期望分解需要独立性及相应可积条件。
对无限求和、积分与期望的交换分别列明非负性或绝对可积性等条件。

\section{方差、矩与位置特征}
\label{sec:ps-variance-moments}
定义\emph{方差}（variance）、标准差、原点矩与中心矩，说明方差有限的二阶矩条件。
推导平移、缩放与方差计算公式；比较均值、中位数和众数的含义。
偏度和峰度作为选读内容，分别注明所需矩存在以及峰度是否减去 $3$ 的约定。

\section{协方差、相关与随机向量的数字特征}
\label{sec:ps-covariance-correlation}
在二阶矩有限条件下定义\emph{协方差}（covariance）与\emph{相关系数}（correlation coefficient），相关系数还要求两个方差均非零。
推导随机变量和的方差，借助柯西--施瓦茨不等式（Cauchy--Schwarz inequality）说明相关系数的范围，并构造不相关但不独立的例子。
引入均值向量与协方差矩阵，证明后者半正定；说明联合正态情形下不相关可推出独立。

\section{例题、习题与学习检查}
规划用指示变量求匹配次数期望、比较有放回和无放回抽样的方差、构造零相关的非线性依赖、检查重尾分布的矩。
统一计算第\ref{chap:ps-univariate}章常见分布的均值和方差，记录参数化及成立条件。
习题要求说明所用性质是否依赖独立性，区分“期望相同”与“分布相同”。






